\section{Harness Engineering：模型如何进入真实工作流}

访谈提到 Harness Engineering，意思是模型本身之外，还需要一层把模型接进真实工作流的工程系统。没有 harness，模型只是一个强接口；有了 harness，模型才能连接代码库、测试、权限、日志、重试、部署和反馈。

\begin{figure}[H]
\centering
\includegraphics[width=0.92\textwidth]{harness-engineering.png}
\caption{Harness Engineering：从模型到真实工作流的中间层。}
\end{figure}

\begin{knowledgebox}{读图：为什么模型强还不够}
图中模型只是第一层。Harness 负责提示、工具、权限、状态和重试；Workflow 是真实代码库和业务流程；Feedback 把测试、用户反馈、成本和失败案例带回系统；Improvement 才能进入后训练和产品迭代。
\end{knowledgebox}

\subsection{术语消化：Harness 相关词表}

\noindent\begin{tabularx}{\textwidth}{p{0.22\textwidth}p{0.32\textwidth}X}
\toprule
术语 & 解决的问题 & 本期中的意义 \\
\midrule
Harness & 把模型接入工具和工作流的外部系统 & 决定模型能否稳定干活。 \\
Workflow & 真实任务流程，如 issue、测试、部署 & 模型价值最终发生在 workflow 中。 \\
Feedback & 测试、用户、成本、失败案例 & 让系统可学习和可改进。 \\
Permission & 权限与安全边界 & 避免模型越权操作。 \\
Retry / Recovery & 出错后重试和恢复 & 决定长程任务能否稳定完成。 \\
\bottomrule
\end{tabularx}

\subsection{为什么它和后训练相关}

Harness 不只是产品工程，也会反过来影响训练。一个好的 harness 能产生更清晰的失败案例、更真实的任务轨迹、更可验证的反馈。这些都能变成后训练数据。EP138 罗福莉谈 Agent post-train 时强调任务环境和 rollout，和这里是同一条逻辑。

\begin{warningbox}{不要把 harness 当 UI 包装}
UI 只是 harness 的一小部分。真正的 harness 包括状态管理、工具执行、安全、日志、评估和反馈回流。只做一个漂亮界面，不能让模型可靠进入生产。
\end{warningbox}

\subsection{本章小结}

Harness Engineering 是模型变成生产力系统的桥。它把模型能力接进真实任务，并把任务反馈带回训练和产品迭代。

